\chapter{Model Serving and Deployment}
\label{ch:model_serving}

\section*{Learning Objectives}

By the end of this chapter, you will be able to:

\begin{itemize}
    \item Design scalable model serving architectures for production ML systems
    \item Implement high-performance REST APIs using FastAPI with proper validation and error handling
    \item Build gRPC services for low-latency model serving
    \item Choose appropriate serving patterns for batch and real-time predictions
    \item Containerize models using Docker and deploy to Kubernetes
    \item Implement safe deployment strategies including blue-green and canary deployments
    \item Optimize serving performance through caching, batching, and model optimization
\end{itemize}

\section{Introduction}

Model serving is the process of making trained models available for prediction in production environments. While training focuses on maximizing model performance, serving emphasizes reliability, latency, throughput, and operational excellence.

\subsection{Serving Architecture Considerations}

\begin{table}[h]
\centering
\caption{Model Serving Design Considerations}
\label{tab:serving_considerations}
\begin{tabular}{|l|p{10cm}|}
\hline
\textbf{Consideration} & \textbf{Description} \\
\hline
Latency & Target response time (ms for real-time, minutes/hours for batch) \\
Throughput & Requests per second the service must handle \\
Availability & Uptime SLA (typically 99.9\% or higher) \\
Scalability & Ability to handle traffic spikes and growth \\
Cost & Infrastructure and operational costs \\
Model Size & Memory and storage requirements \\
Update Frequency & How often models need to be updated \\
\hline
\end{tabular}
\end{table}

\subsection{Serving Pattern Comparison}

\begin{figure}[h]
\centering
\begin{tikzpicture}[
    node distance=1.2cm,
    box/.style={rectangle, draw, text width=2.5cm, align=center, minimum height=0.8cm},
    arrow/.style={->, >=stealth}
]
    % Real-time serving
    \node[box] (client1) {Client};
    \node[box, right=of client1] (api1) {API Server};
    \node[box, right=of api1] (model1) {Model};
    \node[above=0.3cm of api1] (label1) {\textbf{Real-time (Synchronous)}};

    \draw[arrow] (client1) -- node[above] {Request} (api1);
    \draw[arrow] (api1) -- (model1);
    \draw[arrow] (model1) -- (api1);
    \draw[arrow] (api1) -- node[above] {Response} (client1);

    % Batch serving
    \node[box, below=2cm of client1] (client2) {Client};
    \node[box, right=of client2] (queue) {Queue};
    \node[box, right=of queue] (batch) {Batch Job};
    \node[box, below=0.8cm of batch] (storage) {Storage};
    \node[above=0.3cm of queue] (label2) {\textbf{Batch (Asynchronous)}};

    \draw[arrow] (client2) -- (queue);
    \draw[arrow] (queue) -- (batch);
    \draw[arrow] (batch) -- (storage);
    \draw[arrow] (storage) -| (client2);
\end{tikzpicture}
\caption{Real-time vs Batch Serving Patterns}
\label{fig:serving_patterns}
\end{figure}

\section{Model Serving Architectures}

\subsection{Serving Architecture Patterns}

\begin{table}[h]
\centering
\caption{Common Serving Architecture Patterns}
\label{tab:serving_architectures}
\begin{tabular}{|l|l|l|l|}
\hline
\textbf{Pattern} & \textbf{Use Case} & \textbf{Latency} & \textbf{Complexity} \\
\hline
Model-in-Service & Simple models, low traffic & Low & Low \\
Model Service & Dedicated model server & Low & Medium \\
Model as Data & Dynamic model loading & Medium & High \\
Model Ensemble & Multiple models, A/B testing & Medium & High \\
\hline
\end{tabular}
\end{table}

\section{REST API Development with FastAPI}

FastAPI is a modern, high-performance web framework for building APIs with Python, featuring automatic validation, serialization, and documentation.

\subsection{Production FastAPI Model Server}

\begin{lstlisting}[style=python, caption={Production-Grade FastAPI Model Serving}]
from typing import List, Dict, Optional, Any
from fastapi import FastAPI, HTTPException, BackgroundTasks, Request, Depends
from fastapi.responses import JSONResponse
from fastapi.middleware.cors import CORSMiddleware
from pydantic import BaseModel, Field, validator
import numpy as np
import joblib
import logging
import time
from datetime import datetime
from prometheus_client import Counter, Histogram, generate_latest, REGISTRY
import uvicorn

# Metrics
REQUEST_COUNT = Counter('model_requests_total', 'Total model prediction requests', ['model_version', 'status'])
REQUEST_LATENCY = Histogram('model_request_duration_seconds', 'Model request latency', ['model_version'])
PREDICTION_COUNT = Counter('predictions_total', 'Total predictions made', ['model_version'])

# Request/Response models
class PredictionRequest(BaseModel):
    """Request schema for predictions."""
    features: Dict[str, float] = Field(..., description="Feature dictionary")
    model_version: Optional[str] = Field(default="latest", description="Model version to use")

    @validator('features')
    def validate_features(cls, v):
        """Validate features."""
        if not v:
            raise ValueError("Features cannot be empty")
        return v

    class Config:
        schema_extra = {
            "example": {
                "features": {
                    "feature1": 1.0,
                    "feature2": 2.0,
                    "feature3": 3.0
                },
                "model_version": "v1.0"
            }
        }

class PredictionResponse(BaseModel):
    """Response schema for predictions."""
    prediction: float
    probability: Optional[List[float]] = None
    model_version: str
    timestamp: str
    latency_ms: float

class BatchPredictionRequest(BaseModel):
    """Request schema for batch predictions."""
    instances: List[Dict[str, float]]
    model_version: Optional[str] = "latest"

class BatchPredictionResponse(BaseModel):
    """Response schema for batch predictions."""
    predictions: List[float]
    model_version: str
    timestamp: str
    total_latency_ms: float
    num_instances: int

class ModelServer:
    """
    Production model serving class.

    Features:
    - Model loading and caching
    - Version management
    - Preprocessing pipeline
    - Error handling
    - Metrics tracking
    """

    def __init__(self, model_dir: str = "./models"):
        """
        Initialize model server.

        Args:
            model_dir: Directory containing model files
        """
        self.model_dir = model_dir
        self.models: Dict[str, Any] = {}
        self.preprocessors: Dict[str, Any] = {}
        self.logger = logging.getLogger(__name__)

        # Load models
        self.load_models()

    def load_models(self):
        """Load all available model versions."""
        import os
        from pathlib import Path

        model_path = Path(self.model_dir)
        if not model_path.exists():
            self.logger.warning(f"Model directory {self.model_dir} not found")
            return

        # Load all model versions
        for model_file in model_path.glob("model_*.pkl"):
            version = model_file.stem.replace("model_", "")
            try:
                self.models[version] = joblib.load(model_file)
                self.logger.info(f"Loaded model version: {version}")

                # Load corresponding preprocessor if exists
                preprocessor_file = model_path / f"preprocessor_{version}.pkl"
                if preprocessor_file.exists():
                    self.preprocessors[version] = joblib.load(preprocessor_file)
                    self.logger.info(f"Loaded preprocessor for version: {version}")

            except Exception as e:
                self.logger.error(f"Failed to load model {version}: {e}")

        # Set latest version
        if self.models:
            latest_version = sorted(self.models.keys())[-1]
            self.models["latest"] = self.models[latest_version]
            if latest_version in self.preprocessors:
                self.preprocessors["latest"] = self.preprocessors[latest_version]
            self.logger.info(f"Set latest model to version: {latest_version}")

    def preprocess(self, features: Dict[str, float], version: str) -> np.ndarray:
        """
        Preprocess features.

        Args:
            features: Raw features
            version: Model version

        Returns:
            Preprocessed feature array
        """
        # Apply preprocessing if available
        if version in self.preprocessors:
            preprocessor = self.preprocessors[version]
            # Convert dict to DataFrame for preprocessing
            import pandas as pd
            df = pd.DataFrame([features])
            processed = preprocessor.transform(df)
            return processed
        else:
            # Convert dict to array (assuming features are ordered)
            return np.array(list(features.values())).reshape(1, -1)

    def predict(self, features: Dict[str, float], version: str = "latest") -> Dict[str, Any]:
        """
        Make prediction.

        Args:
            features: Input features
            version: Model version

        Returns:
            Prediction result
        """
        start_time = time.time()

        # Check model exists
        if version not in self.models:
            raise ValueError(f"Model version {version} not found")

        try:
            # Preprocess
            X = self.preprocess(features, version)

            # Predict
            model = self.models[version]
            prediction = model.predict(X)[0]

            # Get probability if available
            probability = None
            if hasattr(model, 'predict_proba'):
                probability = model.predict_proba(X)[0].tolist()

            # Calculate latency
            latency_ms = (time.time() - start_time) * 1000

            # Record metrics
            REQUEST_LATENCY.labels(model_version=version).observe(time.time() - start_time)
            PREDICTION_COUNT.labels(model_version=version).inc()

            return {
                'prediction': float(prediction),
                'probability': probability,
                'model_version': version,
                'timestamp': datetime.now().isoformat(),
                'latency_ms': latency_ms
            }

        except Exception as e:
            self.logger.error(f"Prediction failed: {e}")
            REQUEST_COUNT.labels(model_version=version, status='error').inc()
            raise

    def predict_batch(
        self,
        instances: List[Dict[str, float]],
        version: str = "latest"
    ) -> Dict[str, Any]:
        """
        Make batch predictions.

        Args:
            instances: List of feature dictionaries
            version: Model version

        Returns:
            Batch prediction results
        """
        start_time = time.time()

        if version not in self.models:
            raise ValueError(f"Model version {version} not found")

        try:
            # Preprocess all instances
            X_list = [self.preprocess(inst, version) for inst in instances]
            X = np.vstack(X_list)

            # Batch predict
            model = self.models[version]
            predictions = model.predict(X).tolist()

            # Calculate latency
            total_latency_ms = (time.time() - start_time) * 1000

            # Record metrics
            PREDICTION_COUNT.labels(model_version=version).inc(len(instances))

            return {
                'predictions': predictions,
                'model_version': version,
                'timestamp': datetime.now().isoformat(),
                'total_latency_ms': total_latency_ms,
                'num_instances': len(instances)
            }

        except Exception as e:
            self.logger.error(f"Batch prediction failed: {e}")
            REQUEST_COUNT.labels(model_version=version, status='error').inc()
            raise

# Initialize FastAPI app
app = FastAPI(
    title="ML Model Serving API",
    description="Production-grade model serving with FastAPI",
    version="1.0.0"
)

# Add CORS middleware
app.add_middleware(
    CORSMiddleware,
    allow_origins=["*"],
    allow_credentials=True,
    allow_methods=["*"],
    allow_headers=["*"],
)

# Initialize model server
model_server = ModelServer(model_dir="./models")

# Middleware for logging
@app.middleware("http")
async def log_requests(request: Request, call_next):
    """Log all requests."""
    start_time = time.time()
    response = await call_next(request)
    latency = time.time() - start_time

    logging.info(
        f"{request.method} {request.url.path} "
        f"status={response.status_code} latency={latency:.3f}s"
    )

    return response

# Health check endpoint
@app.get("/health")
async def health_check():
    """Health check endpoint."""
    return {
        "status": "healthy",
        "timestamp": datetime.now().isoformat(),
        "models_loaded": list(model_server.models.keys())
    }

# Readiness check endpoint
@app.get("/ready")
async def readiness_check():
    """Readiness check endpoint."""
    if not model_server.models:
        raise HTTPException(status_code=503, detail="No models loaded")

    return {
        "status": "ready",
        "timestamp": datetime.now().isoformat()
    }

# Prediction endpoint
@app.post("/predict", response_model=PredictionResponse)
async def predict(request: PredictionRequest):
    """
    Make a single prediction.

    Args:
        request: Prediction request

    Returns:
        Prediction response
    """
    try:
        result = model_server.predict(
            features=request.features,
            version=request.model_version
        )

        REQUEST_COUNT.labels(
            model_version=request.model_version,
            status='success'
        ).inc()

        return PredictionResponse(**result)

    except ValueError as e:
        raise HTTPException(status_code=400, detail=str(e))
    except Exception as e:
        raise HTTPException(status_code=500, detail=f"Prediction failed: {str(e)}")

# Batch prediction endpoint
@app.post("/predict_batch", response_model=BatchPredictionResponse)
async def predict_batch(request: BatchPredictionRequest):
    """
    Make batch predictions.

    Args:
        request: Batch prediction request

    Returns:
        Batch prediction response
    """
    try:
        result = model_server.predict_batch(
            instances=request.instances,
            version=request.model_version
        )

        REQUEST_COUNT.labels(
            model_version=request.model_version,
            status='success'
        ).inc()

        return BatchPredictionResponse(**result)

    except ValueError as e:
        raise HTTPException(status_code=400, detail=str(e))
    except Exception as e:
        raise HTTPException(status_code=500, detail=f"Batch prediction failed: {str(e)}")

# Model info endpoint
@app.get("/models")
async def list_models():
    """List available model versions."""
    return {
        "models": list(model_server.models.keys()),
        "timestamp": datetime.now().isoformat()
    }

# Metrics endpoint for Prometheus
@app.get("/metrics")
async def metrics():
    """Expose Prometheus metrics."""
    from fastapi.responses import PlainTextResponse
    return PlainTextResponse(generate_latest(REGISTRY))

# Example usage
if __name__ == "__main__":
    # Configure logging
    logging.basicConfig(
        level=logging.INFO,
        format='%(asctime)s - %(name)s - %(levelname)s - %(message)s'
    )

    # Run server
    uvicorn.run(
        app,
        host="0.0.0.0",
        port=8000,
        workers=4,  # Number of worker processes
        log_level="info"
    )
\end{lstlisting}

\subsection{API Testing and Documentation}

FastAPI automatically generates interactive API documentation:

\begin{lstlisting}[style=python, caption={API Testing with pytest}]
import pytest
from fastapi.testclient import TestClient
from main import app, model_server
import numpy as np
from sklearn.ensemble import RandomForestClassifier
import joblib
from pathlib import Path

@pytest.fixture
def client():
    """Create test client."""
    return TestClient(app)

@pytest.fixture(autouse=True)
def setup_test_model(tmp_path):
    """Setup test model."""
    # Create dummy model
    X = np.random.randn(100, 3)
    y = np.random.randint(0, 2, 100)
    model = RandomForestClassifier(n_estimators=10, random_state=42)
    model.fit(X, y)

    # Save model
    model_dir = Path("./models")
    model_dir.mkdir(exist_ok=True)
    joblib.dump(model, model_dir / "model_v1.pkl")

    # Reload models
    model_server.load_models()

    yield

    # Cleanup
    (model_dir / "model_v1.pkl").unlink(missing_ok=True)

def test_health_check(client):
    """Test health check endpoint."""
    response = client.get("/health")
    assert response.status_code == 200
    assert response.json()["status"] == "healthy"

def test_readiness_check(client):
    """Test readiness check endpoint."""
    response = client.get("/ready")
    assert response.status_code == 200
    assert response.json()["status"] == "ready"

def test_predict(client):
    """Test prediction endpoint."""
    request_data = {
        "features": {"feature1": 1.0, "feature2": 2.0, "feature3": 3.0},
        "model_version": "latest"
    }

    response = client.post("/predict", json=request_data)
    assert response.status_code == 200

    data = response.json()
    assert "prediction" in data
    assert "model_version" in data
    assert "latency_ms" in data

def test_predict_batch(client):
    """Test batch prediction endpoint."""
    request_data = {
        "instances": [
            {"feature1": 1.0, "feature2": 2.0, "feature3": 3.0},
            {"feature1": 4.0, "feature2": 5.0, "feature3": 6.0}
        ],
        "model_version": "latest"
    }

    response = client.post("/predict_batch", json=request_data)
    assert response.status_code == 200

    data = response.json()
    assert "predictions" in data
    assert len(data["predictions"]) == 2
    assert data["num_instances"] == 2

def test_invalid_model_version(client):
    """Test with invalid model version."""
    request_data = {
        "features": {"feature1": 1.0, "feature2": 2.0, "feature3": 3.0},
        "model_version": "invalid"
    }

    response = client.post("/predict", json=request_data)
    assert response.status_code == 400

def test_list_models(client):
    """Test listing models."""
    response = client.get("/models")
    assert response.status_code == 200
    assert "models" in response.json()

def test_metrics(client):
    """Test metrics endpoint."""
    response = client.get("/metrics")
    assert response.status_code == 200
    assert "model_requests_total" in response.text
\end{lstlisting}

\section{gRPC for High Performance Serving}

gRPC provides high-performance, low-latency communication using Protocol Buffers for serialization.

\subsection{gRPC Service Implementation}

\begin{lstlisting}[style=protobuf, caption={Protocol Buffer Definition}]
syntax = "proto3";

package ml_serving;

// Prediction service
service PredictionService {
  rpc Predict(PredictionRequest) returns (PredictionResponse);
  rpc PredictStream(stream PredictionRequest) returns (stream PredictionResponse);
  rpc PredictBatch(BatchPredictionRequest) returns (BatchPredictionResponse);
}

// Single prediction request
message PredictionRequest {
  map<string, float> features = 1;
  string model_version = 2;
}

// Single prediction response
message PredictionResponse {
  float prediction = 1;
  repeated float probabilities = 2;
  string model_version = 3;
  int64 timestamp_ms = 4;
  float latency_ms = 5;
}

// Batch prediction request
message BatchPredictionRequest {
  repeated PredictionRequest requests = 1;
  string model_version = 2;
}

// Batch prediction response
message BatchPredictionResponse {
  repeated float predictions = 1;
  string model_version = 2;
  int64 timestamp_ms = 3;
  float total_latency_ms = 4;
  int32 num_instances = 5;
}
\end{lstlisting}

\begin{lstlisting}[style=python, caption={gRPC Server Implementation}]
import grpc
from concurrent import futures
import time
import logging
import numpy as np
import joblib
from typing import Dict, Any

# Import generated protobuf code
# import ml_serving_pb2
# import ml_serving_pb2_grpc

class PredictionServicer:
    """gRPC prediction servicer."""

    def __init__(self, model_path: str):
        """
        Initialize servicer.

        Args:
            model_path: Path to model file
        """
        self.model = joblib.load(model_path)
        self.logger = logging.getLogger(__name__)

    def Predict(self, request, context):
        """Handle single prediction request."""
        start_time = time.time()

        try:
            # Extract features
            features = dict(request.features)
            X = np.array(list(features.values())).reshape(1, -1)

            # Predict
            prediction = self.model.predict(X)[0]

            # Get probabilities if available
            probabilities = []
            if hasattr(self.model, 'predict_proba'):
                probabilities = self.model.predict_proba(X)[0].tolist()

            # Calculate latency
            latency_ms = (time.time() - start_time) * 1000

            # Create response
            response = ml_serving_pb2.PredictionResponse(
                prediction=float(prediction),
                probabilities=probabilities,
                model_version=request.model_version,
                timestamp_ms=int(time.time() * 1000),
                latency_ms=latency_ms
            )

            return response

        except Exception as e:
            self.logger.error(f"Prediction failed: {e}")
            context.set_code(grpc.StatusCode.INTERNAL)
            context.set_details(f"Prediction failed: {str(e)}")
            return ml_serving_pb2.PredictionResponse()

    def PredictStream(self, request_iterator, context):
        """Handle streaming predictions."""
        for request in request_iterator:
            response = self.Predict(request, context)
            yield response

    def PredictBatch(self, request, context):
        """Handle batch prediction request."""
        start_time = time.time()

        try:
            # Extract all features
            X_list = []
            for req in request.requests:
                features = dict(req.features)
                X_list.append(list(features.values()))

            X = np.array(X_list)

            # Batch predict
            predictions = self.model.predict(X).tolist()

            # Calculate latency
            total_latency_ms = (time.time() - start_time) * 1000

            # Create response
            response = ml_serving_pb2.BatchPredictionResponse(
                predictions=predictions,
                model_version=request.model_version,
                timestamp_ms=int(time.time() * 1000),
                total_latency_ms=total_latency_ms,
                num_instances=len(predictions)
            )

            return response

        except Exception as e:
            self.logger.error(f"Batch prediction failed: {e}")
            context.set_code(grpc.StatusCode.INTERNAL)
            context.set_details(f"Batch prediction failed: {str(e)}")
            return ml_serving_pb2.BatchPredictionResponse()

def serve(model_path: str, port: int = 50051, max_workers: int = 10):
    """
    Start gRPC server.

    Args:
        model_path: Path to model file
        port: Server port
        max_workers: Maximum worker threads
    """
    server = grpc.server(futures.ThreadPoolExecutor(max_workers=max_workers))

    # Add servicer
    servicer = PredictionServicer(model_path)
    ml_serving_pb2_grpc.add_PredictionServiceServicer_to_server(servicer, server)

    # Start server
    server.add_insecure_port(f'[::]:{port}')
    server.start()

    logging.info(f"gRPC server started on port {port}")

    try:
        server.wait_for_termination()
    except KeyboardInterrupt:
        server.stop(0)

if __name__ == "__main__":
    logging.basicConfig(level=logging.INFO)
    serve(model_path="./models/model_latest.pkl", port=50051)
\end{lstlisting}

\begin{lstlisting}[style=python, caption={gRPC Client Implementation}]
import grpc
import time
from typing import Dict, List

# import ml_serving_pb2
# import ml_serving_pb2_grpc

class PredictionClient:
    """gRPC prediction client."""

    def __init__(self, host: str = "localhost", port: int = 50051):
        """
        Initialize client.

        Args:
            host: Server host
            port: Server port
        """
        self.channel = grpc.insecure_channel(f'{host}:{port}')
        self.stub = ml_serving_pb2_grpc.PredictionServiceStub(self.channel)

    def predict(
        self,
        features: Dict[str, float],
        model_version: str = "latest"
    ) -> Dict:
        """
        Make single prediction.

        Args:
            features: Feature dictionary
            model_version: Model version

        Returns:
            Prediction result
        """
        request = ml_serving_pb2.PredictionRequest(
            features=features,
            model_version=model_version
        )

        response = self.stub.Predict(request)

        return {
            'prediction': response.prediction,
            'probabilities': list(response.probabilities),
            'model_version': response.model_version,
            'latency_ms': response.latency_ms
        }

    def predict_batch(
        self,
        instances: List[Dict[str, float]],
        model_version: str = "latest"
    ) -> Dict:
        """
        Make batch predictions.

        Args:
            instances: List of feature dictionaries
            model_version: Model version

        Returns:
            Batch prediction results
        """
        requests = [
            ml_serving_pb2.PredictionRequest(features=inst, model_version=model_version)
            for inst in instances
        ]

        batch_request = ml_serving_pb2.BatchPredictionRequest(
            requests=requests,
            model_version=model_version
        )

        response = self.stub.PredictBatch(batch_request)

        return {
            'predictions': list(response.predictions),
            'model_version': response.model_version,
            'total_latency_ms': response.total_latency_ms,
            'num_instances': response.num_instances
        }

    def predict_stream(
        self,
        instances: List[Dict[str, float]],
        model_version: str = "latest"
    ):
        """
        Make streaming predictions.

        Args:
            instances: List of feature dictionaries
            model_version: Model version

        Yields:
            Prediction results
        """
        def request_generator():
            for inst in instances:
                yield ml_serving_pb2.PredictionRequest(
                    features=inst,
                    model_version=model_version
                )

        responses = self.stub.PredictStream(request_generator())

        for response in responses:
            yield {
                'prediction': response.prediction,
                'probabilities': list(response.probabilities),
                'latency_ms': response.latency_ms
            }

    def close(self):
        """Close channel."""
        self.channel.close()

# Example usage
if __name__ == "__main__":
    client = PredictionClient(host="localhost", port=50051)

    # Single prediction
    result = client.predict(
        features={"feature1": 1.0, "feature2": 2.0, "feature3": 3.0}
    )
    print(f"Prediction: {result['prediction']}")
    print(f"Latency: {result['latency_ms']}ms")

    # Batch prediction
    instances = [
        {"feature1": 1.0, "feature2": 2.0, "feature3": 3.0},
        {"feature1": 4.0, "feature2": 5.0, "feature3": 6.0}
    ]
    batch_result = client.predict_batch(instances)
    print(f"\nBatch predictions: {batch_result['predictions']}")
    print(f"Total latency: {batch_result['total_latency_ms']}ms")

    # Streaming predictions
    print("\nStreaming predictions:")
    for result in client.predict_stream(instances):
        print(f"  Prediction: {result['prediction']}, Latency: {result['latency_ms']}ms")

    client.close()
\end{lstlisting}

\section{Batch vs Real-time Serving}

Choosing between batch and real-time serving depends on latency requirements, data availability, and computational costs.

\subsection{Batch Prediction Pipeline}

\begin{lstlisting}[style=python, caption={Scalable Batch Prediction with Apache Beam}]
import apache_beam as beam
from apache_beam.options.pipeline_options import PipelineOptions
import joblib
import numpy as np
import pandas as pd
from typing import Dict, List, Any
import logging

class LoadModel(beam.DoFn):
    """Load model (executed once per worker)."""

    def __init__(self, model_path: str):
        """
        Initialize DoFn.

        Args:
            model_path: Path to model file
        """
        self.model_path = model_path
        self.model = None

    def setup(self):
        """Load model during worker initialization."""
        self.model = joblib.load(self.model_path)
        logging.info(f"Loaded model from {self.model_path}")

    def process(self, element: Dict[str, Any]):
        """
        Make prediction.

        Args:
            element: Feature dictionary

        Yields:
            Prediction result
        """
        # Extract features
        features = {k: v for k, v in element.items() if k != 'id'}
        X = np.array(list(features.values())).reshape(1, -1)

        # Predict
        prediction = self.model.predict(X)[0]

        # Add prediction to element
        element['prediction'] = float(prediction)

        yield element

class BatchPredictor:
    """
    Batch prediction pipeline using Apache Beam.

    Supports local, Dataflow, and other runners.
    """

    def __init__(
        self,
        model_path: str,
        input_pattern: str,
        output_path: str,
        batch_size: int = 1000
    ):
        """
        Initialize batch predictor.

        Args:
            model_path: Path to model file
            input_pattern: Input file pattern (e.g., "gs://bucket/data/*.json")
            output_path: Output file path
            batch_size: Batch size for processing
        """
        self.model_path = model_path
        self.input_pattern = input_pattern
        self.output_path = output_path
        self.batch_size = batch_size

    def run(self, pipeline_options: PipelineOptions = None):
        """
        Run batch prediction pipeline.

        Args:
            pipeline_options: Beam pipeline options
        """
        with beam.Pipeline(options=pipeline_options) as pipeline:
            (
                pipeline
                | 'Read Input' >> beam.io.ReadFromText(self.input_pattern)
                | 'Parse JSON' >> beam.Map(lambda x: eval(x))  # or json.loads
                | 'Batch Elements' >> beam.BatchElements(
                    min_batch_size=self.batch_size,
                    max_batch_size=self.batch_size
                )
                | 'Predict' >> beam.ParDo(LoadModel(self.model_path))
                | 'Format Output' >> beam.Map(lambda x: str(x))
                | 'Write Output' >> beam.io.WriteToText(self.output_path)
            )

# Example usage
if __name__ == "__main__":
    # Configure pipeline
    options = PipelineOptions([
        '--runner=DirectRunner',  # Use DataflowRunner for GCP
        '--project=my-project',
        '--region=us-central1',
        '--temp_location=gs://my-bucket/temp'
    ])

    # Run batch prediction
    predictor = BatchPredictor(
        model_path="./models/model_latest.pkl",
        input_pattern="./data/input/*.json",
        output_path="./data/output/predictions",
        batch_size=1000
    )

    predictor.run(options)
\end{lstlisting}

\subsection{Offline Feature Computation}

\begin{lstlisting}[style=python, caption={Offline Feature Computation for Batch Serving}]
from typing import List, Dict, Any
import pandas as pd
import numpy as np
from datetime import datetime, timedelta
import logging

class OfflineFeatureComputer:
    """
    Compute features offline for batch serving.

    Features are precomputed and stored for fast batch prediction.
    """

    def __init__(self, feature_store_path: str = "./feature_store"):
        """
        Initialize offline feature computer.

        Args:
            feature_store_path: Path to feature store
        """
        self.feature_store_path = feature_store_path
        self.logger = logging.getLogger(__name__)

    def compute_aggregation_features(
        self,
        df: pd.DataFrame,
        entity_column: str,
        timestamp_column: str,
        value_columns: List[str],
        windows: List[int] = [7, 14, 30]
    ) -> pd.DataFrame:
        """
        Compute aggregation features for batch prediction.

        Args:
            df: Input data
            entity_column: Entity identifier column
            timestamp_column: Timestamp column
            value_columns: Columns to aggregate
            windows: Window sizes in days

        Returns:
            DataFrame with aggregation features
        """
        df = df.sort_values([entity_column, timestamp_column])

        for value_col in value_columns:
            for window in windows:
                # Rolling mean
                df[f'{value_col}_rolling_{window}d_mean'] = (
                    df.groupby(entity_column)[value_col]
                    .rolling(window=window, min_periods=1)
                    .mean()
                    .reset_index(level=0, drop=True)
                )

                # Rolling std
                df[f'{value_col}_rolling_{window}d_std'] = (
                    df.groupby(entity_column)[value_col]
                    .rolling(window=window, min_periods=1)
                    .std()
                    .reset_index(level=0, drop=True)
                )

        return df

    def save_features(self, df: pd.DataFrame, feature_set_name: str):
        """
        Save computed features.

        Args:
            df: DataFrame with features
            feature_set_name: Name of feature set
        """
        from pathlib import Path

        path = Path(self.feature_store_path)
        path.mkdir(parents=True, exist_ok=True)

        output_path = path / f"{feature_set_name}.parquet"
        df.to_parquet(output_path, compression='snappy')

        self.logger.info(f"Saved {len(df)} rows to {output_path}")

    def load_features(self, feature_set_name: str) -> pd.DataFrame:
        """
        Load precomputed features.

        Args:
            feature_set_name: Name of feature set

        Returns:
            DataFrame with features
        """
        from pathlib import Path

        path = Path(self.feature_store_path) / f"{feature_set_name}.parquet"
        df = pd.read_parquet(path)

        self.logger.info(f"Loaded {len(df)} rows from {path}")

        return df

# Example usage
if __name__ == "__main__":
    # Generate sample data
    dates = pd.date_range('2023-01-01', periods=365, freq='D')
    entities = ['A', 'B', 'C']

    df = pd.DataFrame({
        'date': np.repeat(dates, len(entities)),
        'entity_id': np.tile(entities, len(dates)),
        'value': np.random.exponential(100, len(dates) * len(entities))
    })

    # Compute features
    computer = OfflineFeatureComputer()
    df_features = computer.compute_aggregation_features(
        df,
        entity_column='entity_id',
        timestamp_column='date',
        value_columns=['value'],
        windows=[7, 14, 30]
    )

    # Save features
    computer.save_features(df_features, "entity_features_2023")

    # Load features
    df_loaded = computer.load_features("entity_features_2023")
    print(f"Loaded features shape: {df_loaded.shape}")
\end{lstlisting}

\section{Model Packaging and Containerization}

Containerization ensures consistent deployment across environments and simplifies dependency management.

\subsection{Docker Containerization}

\begin{lstlisting}[language=docker, caption={Production Dockerfile for Model Serving}]
# Multi-stage build for smaller image size
FROM python:3.9-slim as builder

# Install build dependencies
RUN apt-get update && apt-get install -y \
    gcc \
    g++ \
    && rm -rf /var/lib/apt/lists/*

# Install Python dependencies
COPY requirements.txt .
RUN pip install --user --no-cache-dir -r requirements.txt

# Production image
FROM python:3.9-slim

# Install runtime dependencies
RUN apt-get update && apt-get install -y \
    libgomp1 \
    && rm -rf /var/lib/apt/lists/*

# Create non-root user
RUN useradd -m -u 1000 appuser

# Copy Python dependencies from builder
COPY --from=builder /root/.local /home/appuser/.local

# Set working directory
WORKDIR /app

# Copy application code
COPY --chown=appuser:appuser . .

# Switch to non-root user
USER appuser

# Add local bin to PATH
ENV PATH=/home/appuser/.local/bin:$PATH

# Expose port
EXPOSE 8000

# Health check
HEALTHCHECK --interval=30s --timeout=3s --start-period=5s --retries=3 \
    CMD python -c "import requests; requests.get('http://localhost:8000/health')"

# Run application
CMD ["uvicorn", "main:app", "--host", "0.0.0.0", "--port", "8000", "--workers", "4"]
\end{lstlisting}

\begin{lstlisting}[language=yaml, caption={Docker Compose for Local Development}]
version: '3.8'

services:
  model-server:
    build:
      context: .
      dockerfile: Dockerfile
    ports:
      - "8000:8000"
    volumes:
      - ./models:/app/models:ro  # Mount models as read-only
      - ./logs:/app/logs
    environment:
      - MODEL_DIR=/app/models
      - LOG_LEVEL=INFO
      - WORKERS=4
    restart: unless-stopped
    healthcheck:
      test: ["CMD", "curl", "-f", "http://localhost:8000/health"]
      interval: 30s
      timeout: 3s
      retries: 3

  prometheus:
    image: prom/prometheus:latest
    ports:
      - "9090:9090"
    volumes:
      - ./prometheus.yml:/etc/prometheus/prometheus.yml:ro
      - prometheus-data:/prometheus
    command:
      - '--config.file=/etc/prometheus/prometheus.yml'
      - '--storage.tsdb.path=/prometheus'
    restart: unless-stopped

  grafana:
    image: grafana/grafana:latest
    ports:
      - "3000:3000"
    volumes:
      - grafana-data:/var/lib/grafana
      - ./grafana/dashboards:/etc/grafana/provisioning/dashboards
    environment:
      - GF_SECURITY_ADMIN_PASSWORD=admin
      - GF_USERS_ALLOW_SIGN_UP=false
    restart: unless-stopped

volumes:
  prometheus-data:
  grafana-data:
\end{lstlisting}

\subsection{Kubernetes Deployment}

\begin{lstlisting}[language=yaml, caption={Kubernetes Deployment Manifest}]
apiVersion: apps/v1
kind: Deployment
metadata:
  name: model-server
  labels:
    app: model-server
    version: v1
spec:
  replicas: 3
  selector:
    matchLabels:
      app: model-server
  template:
    metadata:
      labels:
        app: model-server
        version: v1
    spec:
      containers:
      - name: model-server
        image: gcr.io/my-project/model-server:v1.0.0
        ports:
        - containerPort: 8000
          name: http
        env:
        - name: MODEL_DIR
          value: "/models"
        - name: WORKERS
          value: "4"
        resources:
          requests:
            memory: "2Gi"
            cpu: "1000m"
          limits:
            memory: "4Gi"
            cpu: "2000m"
        livenessProbe:
          httpGet:
            path: /health
            port: 8000
          initialDelaySeconds: 30
          periodSeconds: 10
          timeoutSeconds: 3
          failureThreshold: 3
        readinessProbe:
          httpGet:
            path: /ready
            port: 8000
          initialDelaySeconds: 10
          periodSeconds: 5
          timeoutSeconds: 3
          failureThreshold: 3
        volumeMounts:
        - name: model-storage
          mountPath: /models
          readOnly: true
      volumes:
      - name: model-storage
        persistentVolumeClaim:
          claimName: model-pvc
---
apiVersion: v1
kind: Service
metadata:
  name: model-server-service
spec:
  type: LoadBalancer
  selector:
    app: model-server
  ports:
  - port: 80
    targetPort: 8000
    name: http
---
apiVersion: autoscaling/v2
kind: HorizontalPodAutoscaler
metadata:
  name: model-server-hpa
spec:
  scaleTargetRef:
    apiVersion: apps/v1
    kind: Deployment
    name: model-server
  minReplicas: 3
  maxReplicas: 10
  metrics:
  - type: Resource
    resource:
      name: cpu
      target:
        type: Utilization
        averageUtilization: 70
  - type: Resource
    resource:
      name: memory
      target:
        type: Utilization
        averageUtilization: 80
\end{lstlisting}

\section{Deployment Strategies}

Safe deployment strategies minimize risk when rolling out new model versions.

\subsection{Blue-Green Deployment}

\begin{lstlisting}[style=python, caption={Blue-Green Deployment Controller}]
from typing import Dict, Optional
import requests
import time
import logging

class BlueGreenDeployment:
    """
    Blue-Green deployment controller.

    Maintains two identical environments (blue and green).
    Traffic is switched atomically between environments.
    """

    def __init__(
        self,
        blue_endpoint: str,
        green_endpoint: str,
        load_balancer_url: str
    ):
        """
        Initialize blue-green deployment.

        Args:
            blue_endpoint: Blue environment endpoint
            green_endpoint: Green environment endpoint
            load_balancer_url: Load balancer configuration URL
        """
        self.blue_endpoint = blue_endpoint
        self.green_endpoint = green_endpoint
        self.load_balancer_url = load_balancer_url
        self.current_live = "blue"  # Currently serving traffic
        self.logger = logging.getLogger(__name__)

    def deploy_to_green(self, new_model_version: str):
        """
        Deploy new model version to green environment.

        Args:
            new_model_version: New model version to deploy
        """
        self.logger.info(f"Deploying {new_model_version} to green environment")

        # Deploy to green (implementation depends on infrastructure)
        # This could be updating a Kubernetes deployment, ECS task, etc.

        # Wait for deployment to complete
        self._wait_for_healthy(self.green_endpoint)

        self.logger.info("Green environment is healthy")

    def validate_green(self, test_requests: list) -> bool:
        """
        Validate green environment with test requests.

        Args:
            test_requests: List of test requests

        Returns:
            True if validation passes
        """
        self.logger.info("Validating green environment")

        for i, request in enumerate(test_requests):
            try:
                response = requests.post(
                    f"{self.green_endpoint}/predict",
                    json=request,
                    timeout=5
                )

                if response.status_code != 200:
                    self.logger.error(f"Validation failed for request {i}: {response.status_code}")
                    return False

            except Exception as e:
                self.logger.error(f"Validation failed for request {i}: {e}")
                return False

        self.logger.info("Green environment validation passed")
        return True

    def switch_traffic_to_green(self):
        """Switch all traffic from blue to green."""
        self.logger.info("Switching traffic to green environment")

        # Update load balancer to point to green
        # Implementation depends on load balancer (nginx, ALB, etc.)
        self._update_load_balancer("green")

        self.current_live = "green"

        self.logger.info("Traffic switched to green")

    def rollback_to_blue(self):
        """Rollback traffic to blue environment."""
        self.logger.info("Rolling back to blue environment")

        self._update_load_balancer("blue")
        self.current_live = "blue"

        self.logger.info("Rolled back to blue")

    def _wait_for_healthy(self, endpoint: str, timeout: int = 300):
        """
        Wait for endpoint to become healthy.

        Args:
            endpoint: Endpoint to check
            timeout: Maximum wait time in seconds
        """
        start_time = time.time()

        while time.time() - start_time < timeout:
            try:
                response = requests.get(f"{endpoint}/health", timeout=3)
                if response.status_code == 200:
                    return
            except:
                pass

            time.sleep(5)

        raise TimeoutError(f"Endpoint {endpoint} did not become healthy within {timeout}s")

    def _update_load_balancer(self, target: str):
        """
        Update load balancer configuration.

        Args:
            target: Target environment ('blue' or 'green')
        """
        # Implementation depends on load balancer
        # Example: Update nginx config, AWS ALB target group, etc.
        pass

# Example usage
if __name__ == "__main__":
    deployment = BlueGreenDeployment(
        blue_endpoint="http://blue.example.com",
        green_endpoint="http://green.example.com",
        load_balancer_url="http://lb.example.com/config"
    )

    # Deploy to green
    deployment.deploy_to_green("v2.0.0")

    # Validate green
    test_requests = [
        {"features": {"f1": 1.0, "f2": 2.0}},
        {"features": {"f1": 3.0, "f2": 4.0}}
    ]

    if deployment.validate_green(test_requests):
        # Switch traffic
        deployment.switch_traffic_to_green()
    else:
        # Keep blue active
        logging.error("Validation failed, keeping blue environment active")
\end{lstlisting}

\subsection{Canary Deployment}

\begin{lstlisting}[style=python, caption={Canary Deployment Controller}]
from typing import Dict, Optional
import requests
import time
import logging
from dataclasses import dataclass

@dataclass
class CanaryMetrics:
    """Metrics for canary evaluation."""
    error_rate: float
    p50_latency: float
    p95_latency: float
    p99_latency: float
    request_count: int

class CanaryDeployment:
    """
    Canary deployment controller.

    Gradually rolls out new version to subset of traffic,
    monitors metrics, and promotes or rolls back based on health.
    """

    def __init__(
        self,
        stable_endpoint: str,
        canary_endpoint: str,
        prometheus_url: str,
        error_rate_threshold: float = 0.05,
        latency_threshold_ms: float = 500
    ):
        """
        Initialize canary deployment.

        Args:
            stable_endpoint: Stable version endpoint
            canary_endpoint: Canary version endpoint
            prometheus_url: Prometheus URL for metrics
            error_rate_threshold: Maximum acceptable error rate
            latency_threshold_ms: Maximum acceptable P95 latency
        """
        self.stable_endpoint = stable_endpoint
        self.canary_endpoint = canary_endpoint
        self.prometheus_url = prometheus_url
        self.error_rate_threshold = error_rate_threshold
        self.latency_threshold_ms = latency_threshold_ms
        self.logger = logging.getLogger(__name__)

    def deploy_canary(self, canary_traffic_percent: int = 10):
        """
        Deploy canary with specified traffic percentage.

        Args:
            canary_traffic_percent: Percentage of traffic to route to canary
        """
        self.logger.info(f"Deploying canary with {canary_traffic_percent}% traffic")

        # Update traffic routing (implementation depends on infrastructure)
        # Example: Update Istio VirtualService, AWS ALB weighted target groups, etc.
        self._update_traffic_split(canary_traffic_percent)

        self.logger.info("Canary deployed")

    def monitor_canary(self, duration_minutes: int = 10) -> bool:
        """
        Monitor canary health.

        Args:
            duration_minutes: How long to monitor

        Returns:
            True if canary is healthy
        """
        self.logger.info(f"Monitoring canary for {duration_minutes} minutes")

        end_time = time.time() + duration_minutes * 60

        while time.time() < end_time:
            # Get metrics
            stable_metrics = self._get_metrics("stable")
            canary_metrics = self._get_metrics("canary")

            # Check error rate
            if canary_metrics.error_rate > self.error_rate_threshold:
                self.logger.error(
                    f"Canary error rate too high: {canary_metrics.error_rate:.3f} > {self.error_rate_threshold}"
                )
                return False

            # Compare error rate with stable
            if canary_metrics.error_rate > stable_metrics.error_rate * 2:
                self.logger.error(
                    f"Canary error rate significantly higher than stable: "
                    f"{canary_metrics.error_rate:.3f} vs {stable_metrics.error_rate:.3f}"
                )
                return False

            # Check latency
            if canary_metrics.p95_latency > self.latency_threshold_ms:
                self.logger.error(
                    f"Canary latency too high: {canary_metrics.p95_latency:.0f}ms > {self.latency_threshold_ms}ms"
                )
                return False

            # Compare latency with stable
            if canary_metrics.p95_latency > stable_metrics.p95_latency * 1.5:
                self.logger.error(
                    f"Canary latency significantly higher than stable: "
                    f"{canary_metrics.p95_latency:.0f}ms vs {stable_metrics.p95_latency:.0f}ms"
                )
                return False

            self.logger.info(
                f"Canary healthy - Error rate: {canary_metrics.error_rate:.3f}, "
                f"P95 latency: {canary_metrics.p95_latency:.0f}ms"
            )

            time.sleep(30)  # Check every 30 seconds

        return True

    def promote_canary(self):
        """Promote canary to 100% traffic."""
        self.logger.info("Promoting canary to 100% traffic")

        # Gradually increase traffic
        for traffic_percent in [25, 50, 75, 100]:
            self._update_traffic_split(traffic_percent)
            time.sleep(60)  # Wait between increases

            # Monitor during ramp
            if traffic_percent < 100:
                if not self.monitor_canary(duration_minutes=5):
                    self.logger.error("Canary unhealthy during promotion, rolling back")
                    self.rollback_canary()
                    return False

        self.logger.info("Canary successfully promoted")
        return True

    def rollback_canary(self):
        """Rollback canary deployment."""
        self.logger.info("Rolling back canary")

        self._update_traffic_split(0)  # Route 0% traffic to canary

        self.logger.info("Canary rolled back")

    def _get_metrics(self, version: str) -> CanaryMetrics:
        """
        Get metrics from Prometheus.

        Args:
            version: Version to get metrics for ('stable' or 'canary')

        Returns:
            Canary metrics
        """
        # Query Prometheus
        # This is a simplified example - actual implementation would use promql
        try:
            # Error rate query
            error_rate_query = f'rate(http_requests_total{{version="{version}",status=~"5.."}}[5m])'
            error_rate = self._query_prometheus(error_rate_query)

            # Latency queries
            p50_query = f'histogram_quantile(0.50, rate(http_request_duration_seconds_bucket{{version="{version}"}}[5m]))'
            p95_query = f'histogram_quantile(0.95, rate(http_request_duration_seconds_bucket{{version="{version}"}}[5m]))'
            p99_query = f'histogram_quantile(0.99, rate(http_request_duration_seconds_bucket{{version="{version}"}}[5m]))'

            p50_latency = self._query_prometheus(p50_query) * 1000  # Convert to ms
            p95_latency = self._query_prometheus(p95_query) * 1000
            p99_latency = self._query_prometheus(p99_query) * 1000

            # Request count
            count_query = f'sum(rate(http_requests_total{{version="{version}"}}[5m]))'
            request_count = self._query_prometheus(count_query)

            return CanaryMetrics(
                error_rate=error_rate,
                p50_latency=p50_latency,
                p95_latency=p95_latency,
                p99_latency=p99_latency,
                request_count=int(request_count)
            )

        except Exception as e:
            self.logger.error(f"Failed to get metrics: {e}")
            # Return safe defaults that will trigger rollback
            return CanaryMetrics(
                error_rate=1.0,
                p50_latency=10000,
                p95_latency=10000,
                p99_latency=10000,
                request_count=0
            )

    def _query_prometheus(self, query: str) -> float:
        """Query Prometheus and return result."""
        try:
            response = requests.get(
                f"{self.prometheus_url}/api/v1/query",
                params={'query': query},
                timeout=5
            )
            data = response.json()
            if data['status'] == 'success' and data['data']['result']:
                return float(data['data']['result'][0]['value'][1])
            return 0.0
        except:
            return 0.0

    def _update_traffic_split(self, canary_percent: int):
        """
        Update traffic split between stable and canary.

        Args:
            canary_percent: Percentage to route to canary
        """
        # Implementation depends on infrastructure
        # Example: Update Istio, AWS ALB, nginx, etc.
        pass

# Example usage
if __name__ == "__main__":
    deployment = CanaryDeployment(
        stable_endpoint="http://stable.example.com",
        canary_endpoint="http://canary.example.com",
        prometheus_url="http://prometheus.example.com",
        error_rate_threshold=0.05,
        latency_threshold_ms=500
    )

    # Deploy canary with 10% traffic
    deployment.deploy_canary(canary_traffic_percent=10)

    # Monitor canary
    if deployment.monitor_canary(duration_minutes=10):
        # Promote if healthy
        deployment.promote_canary()
    else:
        # Rollback if unhealthy
        deployment.rollback_canary()
\end{lstlisting}

\section{Serving Optimization and Caching}

Optimizing serving performance reduces latency and costs through caching, batching, and model optimization.

\subsection{Response Caching}

\begin{lstlisting}[style=python, caption={Intelligent Response Caching}]
from typing import Dict, Any, Optional, Callable
import hashlib
import json
import time
from functools import wraps
import redis
import logging

class PredictionCache:
    """
    Intelligent caching for model predictions.

    Features:
    - Redis-based distributed caching
    - TTL support
    - Cache hit/miss tracking
    - Automatic cache warming
    """

    def __init__(
        self,
        redis_host: str = "localhost",
        redis_port: int = 6379,
        default_ttl: int = 3600,  # 1 hour
        enable_metrics: bool = True
    ):
        """
        Initialize prediction cache.

        Args:
            redis_host: Redis host
            redis_port: Redis port
            default_ttl: Default TTL in seconds
            enable_metrics: Enable cache metrics tracking
        """
        self.redis_client = redis.Redis(
            host=redis_host,
            port=redis_port,
            decode_responses=True
        )
        self.default_ttl = default_ttl
        self.enable_metrics = enable_metrics
        self.logger = logging.getLogger(__name__)

        # Metrics
        self.cache_hits = 0
        self.cache_misses = 0

    def _generate_cache_key(self, features: Dict[str, Any], model_version: str) -> str:
        """
        Generate deterministic cache key.

        Args:
            features: Feature dictionary
            model_version: Model version

        Returns:
            Cache key
        """
        # Sort features for deterministic key
        sorted_features = json.dumps(features, sort_keys=True)
        key_string = f"{model_version}:{sorted_features}"

        # Hash for fixed-length key
        return hashlib.sha256(key_string.encode()).hexdigest()

    def get(self, features: Dict[str, Any], model_version: str) -> Optional[Dict[str, Any]]:
        """
        Get prediction from cache.

        Args:
            features: Feature dictionary
            model_version: Model version

        Returns:
            Cached prediction or None
        """
        cache_key = self._generate_cache_key(features, model_version)

        try:
            cached_value = self.redis_client.get(cache_key)

            if cached_value:
                if self.enable_metrics:
                    self.cache_hits += 1

                self.logger.debug(f"Cache hit for key {cache_key}")
                return json.loads(cached_value)
            else:
                if self.enable_metrics:
                    self.cache_misses += 1

                self.logger.debug(f"Cache miss for key {cache_key}")
                return None

        except Exception as e:
            self.logger.error(f"Cache get failed: {e}")
            return None

    def set(
        self,
        features: Dict[str, Any],
        model_version: str,
        prediction: Dict[str, Any],
        ttl: Optional[int] = None
    ):
        """
        Store prediction in cache.

        Args:
            features: Feature dictionary
            model_version: Model version
            prediction: Prediction result
            ttl: Time-to-live in seconds
        """
        cache_key = self._generate_cache_key(features, model_version)
        ttl = ttl or self.default_ttl

        try:
            self.redis_client.setex(
                cache_key,
                ttl,
                json.dumps(prediction)
            )

            self.logger.debug(f"Cached prediction for key {cache_key}")

        except Exception as e:
            self.logger.error(f"Cache set failed: {e}")

    def invalidate(self, model_version: Optional[str] = None):
        """
        Invalidate cache entries.

        Args:
            model_version: Specific version to invalidate (None = all)
        """
        if model_version:
            # Invalidate specific version (requires key pattern tracking)
            pattern = f"*{model_version}*"
            keys = self.redis_client.keys(pattern)
            if keys:
                self.redis_client.delete(*keys)
                self.logger.info(f"Invalidated {len(keys)} cache entries for {model_version}")
        else:
            # Invalidate all
            self.redis_client.flushdb()
            self.logger.info("Invalidated all cache entries")

    def get_metrics(self) -> Dict[str, Any]:
        """Get cache metrics."""
        total_requests = self.cache_hits + self.cache_misses
        hit_rate = self.cache_hits / total_requests if total_requests > 0 else 0

        return {
            'cache_hits': self.cache_hits,
            'cache_misses': self.cache_misses,
            'total_requests': total_requests,
            'hit_rate': hit_rate
        }

def cached_prediction(cache: PredictionCache, ttl: Optional[int] = None):
    """
    Decorator for caching predictions.

    Args:
        cache: PredictionCache instance
        ttl: Time-to-live for cache entry

    Returns:
        Decorated function
    """
    def decorator(predict_func: Callable):
        @wraps(predict_func)
        def wrapper(features: Dict[str, Any], model_version: str = "latest", **kwargs):
            # Try cache first
            cached_result = cache.get(features, model_version)
            if cached_result is not None:
                cached_result['from_cache'] = True
                return cached_result

            # Call actual prediction
            result = predict_func(features, model_version, **kwargs)

            # Cache result
            cache.set(features, model_version, result, ttl=ttl)

            result['from_cache'] = False
            return result

        return wrapper

    return decorator

# Example usage
if __name__ == "__main__":
    # Initialize cache
    cache = PredictionCache(redis_host="localhost", redis_port=6379)

    # Example prediction function
    @cached_prediction(cache, ttl=3600)
    def predict(features: Dict[str, Any], model_version: str = "latest"):
        # Simulate expensive prediction
        time.sleep(0.1)

        return {
            'prediction': 0.8,
            'model_version': model_version,
            'timestamp': time.time()
        }

    # First call - cache miss
    features = {"feature1": 1.0, "feature2": 2.0}
    result1 = predict(features)
    print(f"First call - From cache: {result1.get('from_cache')}")

    # Second call - cache hit
    result2 = predict(features)
    print(f"Second call - From cache: {result2.get('from_cache')}")

    # Get metrics
    metrics = cache.get_metrics()
    print(f"\nCache metrics: {metrics}")
\end{lstlisting}

\subsection{Request Batching}

\begin{lstlisting}[style=python, caption={Dynamic Request Batching}]
from typing import List, Dict, Any, Callable
import asyncio
import time
from dataclasses import dataclass
import logging

@dataclass
class BatchConfig:
    """Configuration for request batching."""
    max_batch_size: int = 32
    max_wait_time_ms: float = 10.0  # Maximum time to wait for batch
    timeout_seconds: float = 30.0

class RequestBatcher:
    """
    Dynamic request batching for efficient inference.

    Collects requests and processes them in batches to maximize throughput.
    """

    def __init__(
        self,
        predict_func: Callable,
        config: BatchConfig = None
    ):
        """
        Initialize request batcher.

        Args:
            predict_func: Batch prediction function
            config: Batch configuration
        """
        self.predict_func = predict_func
        self.config = config or BatchConfig()

        self.pending_requests: List[Dict[str, Any]] = []
        self.pending_futures: List[asyncio.Future] = []
        self.batch_lock = asyncio.Lock()
        self.logger = logging.getLogger(__name__)

        # Start batch processor
        asyncio.create_task(self._process_batches())

    async def predict(self, features: Dict[str, Any]) -> Dict[str, Any]:
        """
        Add prediction request to batch.

        Args:
            features: Feature dictionary

        Returns:
            Prediction result
        """
        # Create future for this request
        future = asyncio.Future()

        async with self.batch_lock:
            self.pending_requests.append(features)
            self.pending_futures.append(future)

            # Trigger immediate batch if max size reached
            if len(self.pending_requests) >= self.config.max_batch_size:
                asyncio.create_task(self._execute_batch())

        # Wait for result
        try:
            result = await asyncio.wait_for(future, timeout=self.config.timeout_seconds)
            return result
        except asyncio.TimeoutError:
            self.logger.error("Request timeout")
            raise

    async def _process_batches(self):
        """Background task to process batches periodically."""
        while True:
            await asyncio.sleep(self.config.max_wait_time_ms / 1000)

            async with self.batch_lock:
                if self.pending_requests:
                    asyncio.create_task(self._execute_batch())

    async def _execute_batch(self):
        """Execute batch prediction."""
        async with self.batch_lock:
            if not self.pending_requests:
                return

            # Get current batch
            batch_requests = self.pending_requests.copy()
            batch_futures = self.pending_futures.copy()

            # Clear pending
            self.pending_requests.clear()
            self.pending_futures.clear()

        self.logger.info(f"Processing batch of {len(batch_requests)} requests")

        try:
            # Execute batch prediction
            start_time = time.time()
            predictions = await self._run_batch_prediction(batch_requests)
            batch_time = (time.time() - start_time) * 1000

            self.logger.info(f"Batch processed in {batch_time:.2f}ms")

            # Distribute results to futures
            for future, prediction in zip(batch_futures, predictions):
                if not future.done():
                    future.set_result(prediction)

        except Exception as e:
            self.logger.error(f"Batch prediction failed: {e}")

            # Set exception for all futures
            for future in batch_futures:
                if not future.done():
                    future.set_exception(e)

    async def _run_batch_prediction(self, batch: List[Dict[str, Any]]) -> List[Dict[str, Any]]:
        """
        Run batch prediction.

        Args:
            batch: List of feature dictionaries

        Returns:
            List of predictions
        """
        # Run prediction in thread pool to avoid blocking
        loop = asyncio.get_event_loop()
        predictions = await loop.run_in_executor(None, self.predict_func, batch)

        return predictions

# Example usage
if __name__ == "__main__":
    import numpy as np
    from sklearn.ensemble import RandomForestClassifier
    import joblib

    # Train dummy model
    X = np.random.randn(100, 3)
    y = np.random.randint(0, 2, 100)
    model = RandomForestClassifier(n_estimators=10)
    model.fit(X, y)

    # Batch prediction function
    def batch_predict(batch: List[Dict[str, Any]]) -> List[Dict[str, Any]]:
        """Process batch of requests."""
        # Convert to array
        X = np.array([[req['f1'], req['f2'], req['f3']] for req in batch])

        # Batch predict
        predictions = model.predict(X)

        # Format results
        results = [
            {'prediction': float(pred), 'from_batch': True}
            for pred in predictions
        ]

        return results

    # Initialize batcher
    batcher = RequestBatcher(
        predict_func=batch_predict,
        config=BatchConfig(max_batch_size=32, max_wait_time_ms=10)
    )

    # Simulate requests
    async def simulate_requests():
        tasks = []
        for i in range(100):
            features = {'f1': float(i), 'f2': float(i+1), 'f3': float(i+2)}
            tasks.append(batcher.predict(features))

        results = await asyncio.gather(*tasks)

        print(f"Processed {len(results)} requests")
        print(f"First result: {results[0]}")

    # Run simulation
    asyncio.run(simulate_requests())
\end{lstlisting}

\section{Chapter Summary}

This chapter presented comprehensive approaches to model serving and deployment:

\begin{itemize}
    \item \textbf{REST API Serving:} Production FastAPI implementation with validation, error handling, and Prometheus metrics

    \item \textbf{gRPC Serving:} High-performance serving using Protocol Buffers for low-latency applications

    \item \textbf{Batch vs Real-time:} Scalable batch prediction with Apache Beam and offline feature computation

    \item \textbf{Containerization:} Docker multi-stage builds and Kubernetes deployments with health checks and auto-scaling

    \item \textbf{Deployment Strategies:} Blue-green deployments for atomic switches and canary deployments with automated rollback

    \item \textbf{Serving Optimization:} Response caching with Redis and dynamic request batching for improved throughput
\end{itemize}

\subsection{Key Takeaways}

\begin{enumerate}
    \item \textbf{Choose the Right Pattern:} Real-time for low-latency requirements, batch for cost-effective bulk processing

    \item \textbf{Validate Inputs:} Use Pydantic or similar for request validation to prevent errors and security issues

    \item \textbf{Implement Health Checks:} Proper liveness and readiness probes for reliable orchestration

    \item \textbf{Monitor Everything:} Track latency, error rates, and throughput to detect issues early

    \item \textbf{Deploy Safely:} Use blue-green or canary deployments to minimize risk

    \item \textbf{Optimize Performance:} Caching, batching, and model optimization reduce costs and improve UX

    \item \textbf{Container Best Practices:} Multi-stage builds, non-root users, and health checks for security and reliability

    \item \textbf{Plan for Scale:} Horizontal pod autoscaling and load balancing for handling traffic growth
\end{enumerate}

\section{Further Reading}

\begin{itemize}
    \item Richardson, C. (2018). \textit{Microservices Patterns}
    \item Burns, B., et al. (2019). \textit{Kubernetes: Up and Running}
    \item Kleppmann, M. (2017). \textit{Designing Data-Intensive Applications}
    \item Humble, J., \& Farley, D. (2010). \textit{Continuous Delivery}
    \item Newman, S. (2021). \textit{Building Microservices, 2nd Edition}
\end{itemize}

\section{Exercises}

\begin{enumerate}
    \item \textbf{FastAPI Model Server:}
    \begin{itemize}
        \item Implement a complete FastAPI server with prediction endpoints
        \item Add request validation using Pydantic
        \item Implement health and readiness checks
        \item Add Prometheus metrics
        \item Write comprehensive tests with pytest
    \end{itemize}

    \item \textbf{gRPC Service:}
    \begin{itemize}
        \item Define Protocol Buffer service
        \item Implement gRPC server and client
        \item Add streaming prediction support
        \item Compare performance with REST API
    \end{itemize}

    \item \textbf{Batch Prediction Pipeline:}
    \begin{itemize}
        \item Implement batch prediction using Apache Beam
        \item Process large dataset (millions of rows)
        \item Optimize for cost and throughput
        \item Add error handling and retry logic
    \end{itemize}

    \item \textbf{Docker Containerization:}
    \begin{itemize}
        \item Create multi-stage Dockerfile
        \item Optimize image size
        \item Add health checks
        \item Run security scanning
        \item Deploy to Kubernetes
    \end{itemize}

    \item \textbf{Deployment Strategy:}
    \begin{itemize}
        \item Implement canary deployment controller
        \item Monitor metrics and auto-rollback
        \item Test with simulated traffic
        \item Document rollback procedures
    \end{itemize}

    \item \textbf{Serving Optimization:}
    \begin{itemize}
        \item Implement Redis caching
        \item Add request batching
        \item Measure latency improvements
        \item Analyze cache hit rates
        \item Optimize for cost vs performance
    \end{itemize}

    \item \textbf{End-to-End Deployment:}
    \begin{itemize}
        \item Build complete serving pipeline from training to production
        \item Implement REST and gRPC endpoints
        \item Add caching and batching
        \item Containerize and deploy to Kubernetes
        \item Implement canary deployment
        \item Monitor with Prometheus and Grafana
    \end{itemize}
\end{enumerate}
